\subsection{Restriction of analytic sections of a vector bundle}

Let X be a complex analytic manifold locally of finite dimension and E an analytic vector bundle with base X (VAR, p. 35 and VAR, p. 70). Denote by \(\mathcal{S}^{\omega}(\mathrm{X};\mathrm{E})\) the vector space of analytic sections of E over X, equipped with the topology of compact convergence.

Lemma 5. --- Let \(X_{0}\) be the underlying real analytic manifold of X and \(E_{0}\) the complex vector bundle over \(X_{0}\) underlying E.

\begin{enumerate}
\def\labelenumi{\alph{enumi})}
\item
  The vector subspace \(\mathcal{S}^{\omega}(\mathrm{X};\mathrm{E})\) of \(\mathcal{S}^{0}(\mathrm{X}_{0};\mathrm{E}_{0})\) is closed, and the injection of \(\mathcal{S}^{\omega}(\mathrm{X};\mathrm{E})\) into \(\mathcal{S}^{0}(\mathrm{X}_{0};\mathrm{E}_{0})\) is continuous and strict;
\item
  The canonical injection of \(\mathcal{S}^{\omega}(\mathrm{X};\mathrm{E})\) into \(\mathcal{S}^{1}(\mathrm{X}_{0};\mathrm{E}_{0})\) is continuous.
\end{enumerate}

Suppose first that there exists an integer \(n \geqslant 0\) and a Banach space F such that X is an open subset of \(C^{n}\) and E is the trivial vector bundle \(X \times F\) . In this case, the topological vector spaces \(\mathcal{S}^{\omega}(X;E)\) , \(\mathcal{S}^{0}(X_{0};E_{0})\) and \(\mathcal{S}^{1}(X_{0};E_{0})\) are isomorphic to \(\mathcal{C}^{\omega}(X;F)\) , \(\mathcal{C}^{0}(X_{0};F)\) and \(\mathcal{C}^{1}(X_{0};F)\) respectively, and the space \(\mathcal{S}^{0}(X_{0};E_{0})\) is metrizable (TG, X, p. 20, cor. of prop. 1). The lemma then follows from VAR, 3.3.2, p. 28.

Let \(\mathfrak{F}\) be a filter on \(\mathcal{S}^{\omega}(\mathrm{X};\mathrm{E})\) which converges in the space \(\mathcal{S}^{0}(\mathrm{X}_{0};\mathrm{E}_{0})\) toward a limit \(f\). We must show that \(f\) belongs to \(\mathcal{S}^{\omega}(\mathrm{X};\mathrm{E})\) and that the filter \(\mathfrak{F}\) converges to \(f\) in the space \(\mathcal{S}^{1}(\mathrm{X}_{0};\mathrm{E}_{0})\). This statement is local in nature, and therefore follows from the first part of the proof.

PROPOSITION 11. --- Suppose that the vector bundle E is locally of finite rank. Let Y be a relatively compact open subset of X. The restriction map \(f \mapsto f|Y\) from \(\mathcal{S}^{\omega}(X;E)\) into \(\mathcal{S}^{\omega}(Y;E)\) is compact.

With the notations of Lemma 5, we have a commutative diagram

\[
\begin{array}{c c c} \mathcal {S} ^ {\omega} (\mathrm{X}; \mathrm{E}) & \stackrel {{i}} {{\longrightarrow}} & \mathcal {S} ^ {1} (\mathrm{X} _ {0}; \mathrm{E} _ {0}) \\ \Big \downarrow^ {u} & & \Big \downarrow^ {v} \\ \mathcal {S} ^ {\omega} (\mathrm{Y}; \mathrm{E}) & \stackrel {{j}} {{\longrightarrow}} & \mathcal {S} ^ {0} (\mathrm{Y} _ {0}; \mathrm{E} _ {0}) \end{array}\tag{3}
\]

where i and j are the canonical injections and u, v the restriction maps. The map i is continuous (Lemma 5, b)), and the map v is compact (Prop. 10). Since the canonical injection j is continuous, strict, and has closed image (Lemma 5, a)), the map u is compact (Remark 5 of III, p. 3).

COROLLARY. --- Let X be a compact complex analytic manifold and E an analytic vector bundle over X, locally of finite rank. The vector space \(\mathcal{S}^{\omega}(X;E)\) is of finite dimension.

Being compact, the analytic manifold X is locally of finite dimension. By prop. 11, the identity map of \(\mathcal{S}^{\omega}(\mathrm{X};\mathrm{E})\) is a compact linear map. This implies that the space \(\mathcal{S}^{\omega}(\mathrm{X};\mathrm{E})\) is of finite dimension (Remark 3 of III, p. 2).

